\documentclass[11pt,a4paper]{article}

\usepackage[T1]{fontenc}
\usepackage[utf8]{inputenc}
\usepackage[french]{babel}
\usepackage{lmodern}
\usepackage{amsmath,amssymb}
\usepackage{booktabs}
\usepackage{longtable}
\usepackage{array}
\usepackage{geometry}
\usepackage{xcolor}
\usepackage{listings}
\usepackage{hyperref}
\usepackage{enumitem}
\usepackage{microtype}

\geometry{margin=2.3cm}
\hypersetup{colorlinks=true,linkcolor=blue!50!black,urlcolor=blue!60!black}
\setlist{nosep}
\lstset{
  basicstyle=\ttfamily\small,
  frame=single,
  breaklines=true,
  columns=fullflexible,
  keepspaces=true,
  showstringspaces=false
}

\newcommand{\code}[1]{\texttt{#1}}
\newcommand{\R}{\mathbb{R}}

\title{Documentation technique du projet CMPUF\\
\large Modèle stochastique, données, procédure expérimentale et résultats}
\author{Projet CMPUF}
\date{\today}

\begin{document}
\maketitle
\tableofcontents
\newpage

\section{Objet et périmètre}

Le programme \code{cmpuf} résout un problème de localisation
d'établissements capacitaires soumis à des perturbations aléatoires de
capacité. Dans le code et les commentaires, ce problème est nommé CMPUF ou
MPUF (\emph{Unreliable Capacitated Facility Location Problem}).

L'implantation des établissements constitue la décision de premier niveau.
Elle est commune à tous les scénarios. Après réalisation d'un scénario de
capacité, les quantités expédiées vers les clients sont réoptimisées. Le
programme met donc en œuvre un programme stochastique discret à deux étapes
avec recours continu.

La cible maintenue est construite à partir de :
\begin{itemize}
  \item \code{cpp/cmpuf.cpp} : lecture des données, scénarios, indicateurs et
        programme principal ;
  \item \code{cpp/modele.cpp} : formulation CPLEX, heuristiques et recherche
        locale ;
  \item \code{cpp/pmedian.cpp} : prétraitement p-médiane ;
  \item \code{cpp/transport.cpp} : résolution indépendante d'un scénario ;
  \item \code{cpp/tri.cpp} : tri des voisinages par distance.
\end{itemize}

\section{Notations et données}

\subsection{Ensembles et indices}

\begin{longtable}{>{$}l<{$} p{11cm}}
\toprule
\text{Notation} & \text{Définition}\\
\midrule
\endhead
N=\{0,\dots,n-1\} & Ensemble des nœuds. Chaque nœud peut être un client et un
emplacement candidat.\\
I=\{0,\dots,m-1\} & Ensemble des établissements à implanter.\\
R=\{0,\dots,ncap-1\} & Ensemble des états de capacité possibles d'un
établissement.\\
\Omega=\{0,\dots,K-1\} & Ensemble des scénarios de capacité, avec
$K=ncap^m$.\\
i & Indice d'établissement.\\
l & Indice d'emplacement candidat.\\
j & Indice de client.\\
r & Indice d'état de capacité.\\
k & Indice de scénario.\\
\bottomrule
\end{longtable}

\subsection{Paramètres}

\begin{longtable}{>{$}l<{$} p{11cm}}
\toprule
\text{Paramètre} & \text{Définition}\\
\midrule
\endhead
w_j & Demande du client $j$.\\
H=\sum_j w_j & Demande totale.\\
d_{lj} & Coût ou distance entre l'emplacement $l$ et le client $j$ ; variable
\code{profit[l][j]} dans le code.\\
D_{\max}=\max_{l,j}d_{lj} & Plus grande distance de l'instance.\\
C^m_i & Capacité médiane : demande affectée à l'établissement $i$ dans la
solution du prétraitement p-médiane.\\
\gamma & Facteur appliqué à la capacité médiane.\\
S_i=\gamma C^m_i & Capacité nominale de l'établissement $i$.\\
b_r & Multiplicateur de capacité lu dans \code{wbeta} pour $r<ncap-1$.\\
a_r & Poids probabiliste lu dans \code{walpha} pour $r<ncap-1$.\\
\alpha & Niveau global de perturbation.\\
p_r & Probabilité marginale de l'état $r$.\\
C_{ki} & Capacité de l'établissement $i$ dans le scénario $k$.\\
P_k & Probabilité du scénario $k$.\\
\bottomrule
\end{longtable}

Le membre \code{beta} est actuellement lu, stocké et écrit dans les résultats,
mais n'intervient dans aucune formule active de capacité. Les multiplicateurs
effectifs sont directement les valeurs \code{wbeta}.

\section{Construction des capacités et des scénarios}

\subsection{Capacités nominales et dégradées}

Le p-médian fournit $C^m_i$. La capacité nominale est :
\begin{equation}
S_i=\gamma C^m_i.
\end{equation}

Les états de capacité sont :
\begin{align}
\operatorname{cap}_{ir} &= b_r S_i,
  &&r=0,\dots,ncap-2,\\
\operatorname{cap}_{i,ncap-1} &= S_i.
\end{align}

Avec le fichier \code{parameters/param} fourni, on obtient quatre états :
$0$, $0{,}33S_i$, $0{,}66S_i$ et $S_i$.

\subsection{Probabilités marginales}

Le code calcule :
\begin{align}
p_r &= \alpha a_r,
  &&r=0,\dots,ncap-2,\\
p_{ncap-1} &= 1-\alpha.
\end{align}

Pour que $(p_r)$ soit une loi de probabilité, les poids $a_r$ doivent sommer
à un. C'est le cas du fichier fourni :
$0{,}10+0{,}40+0{,}50=1$.

\subsection{Énumération exhaustive}

Un scénario est un vecteur
\[
T_k=(T_{k0},\dots,T_{k,m-1})\in R^m.
\]
La fonction récursive \code{Init\_T} énumère le produit cartésien complet
$R^m$. Pour chaque scénario :
\begin{align}
C_{ki} &= \operatorname{cap}_{i,T_{ki}},\\
P_k &= \prod_{i\in I}p_{T_{ki}}.
\end{align}

La croissance est exponentielle :
\[
K=ncap^m.
\]
Le modèle principal contient de l'ordre de $n^2K$ variables de flux. Cette
énumération constitue la principale limite de passage à l'échelle.

\section{Modèle mathématique stochastique codé}

\subsection{Variables de décision}

Les variables de premier niveau sont :
\[
x_{il}=
\begin{cases}
1 & \text{si l'établissement $i$ est implanté au nœud $l$,}\\
0 & \text{sinon.}
\end{cases}
\]

Dans Concert, les variables sont initialement continues dans $[0,1]$, puis
converties en variables entières par \code{IloConversion}. Elles sont donc
binaires après conversion.

Les variables de recours sont :
\[
z_{ljk}\geq 0,
\]
quantité expédiée depuis le nœud $l$ vers le client $j$ dans le scénario $k$.

\subsection{Objectif}

Le programme maximise l'utilité espérée :
\begin{equation}
\max\quad
\sum_{k\in\Omega}P_k
\sum_{l\in N}\sum_{j\in N}
(D_{\max}-d_{lj})z_{ljk}.
\label{eq:objectif}
\end{equation}

Lorsque la quantité totale livrée est fixée, maximiser
\eqref{eq:objectif} revient à minimiser le coût de transport espéré.
La transformation maintient des coefficients non négatifs.

\paragraph{Point d'attention.}
Dans le modèle codé, les demandes sont des bornes inférieures et non des
égalités. Si une capacité excédentaire existe et si
$D_{\max}-d_{lj}>0$, le solveur peut augmenter certains flux au-delà de la
demande. Dans ce cas, l'objectif codé n'est plus strictement équivalent à une
minimisation du coût pour une demande exactement satisfaite.

\subsection{Contraintes d'implantation}

Chaque établissement doit être placé exactement une fois :
\begin{equation}
\sum_{l\in N}x_{il}=1,
\qquad i\in I.
\label{eq:localisation}
\end{equation}

Dans le chemin d'exécution actuel, ces contraintes sont ajoutées
progressivement par l'heuristique gloutonne et restent ensuite dans le modèle
utilisé pour la résolution exacte.

La méthode \code{Init\_contraintes\_x\_x} permettrait d'interdire la
colocalisation :
\[
\sum_{i\in I}x_{il}\leq 1,\qquad l\in N,
\]
mais son appel est commenté. Plusieurs établissements peuvent donc occuper le
même nœud.

\subsection{Contraintes de capacité}

Pour chaque emplacement et scénario :
\begin{equation}
\sum_{j\in N}z_{ljk}
\leq
\sum_{i\in I}C_{ki}x_{il},
\qquad l\in N,\ k\in\Omega.
\label{eq:capacite}
\end{equation}

La capacité disponible en un nœud est la somme des capacités des
établissements qui y sont colocalisés.

\subsection{Contraintes de demande}

Pour chaque client et scénario :
\begin{equation}
\sum_{l\in N}z_{ljk}\geq w_j,
\qquad j\in N,\ k\in\Omega.
\label{eq:demande}
\end{equation}

Les variables de flux sont non négatives et aucune borne supérieure directe
par client n'est imposée.

\subsection{Structure à deux étapes}

Les $x_{il}$ sont identiques pour tous les scénarios : ce sont les décisions
\emph{ici et maintenant}. Les $z_{ljk}$ dépendent de $k$ : ils constituent le
recours après observation des capacités. Le code peut annoter les variables
pour la décomposition de Benders :
\begin{itemize}
  \item les $x$ sont placés dans le maître ;
  \item les flux $z$ sont répartis par scénario.
\end{itemize}
La stratégie Benders de la résolution exacte est toutefois réglée à $-1$,
c'est-à-dire laissée au choix automatique de CPLEX.

\section{Prétraitement p-médiane}

Avant le modèle stochastique, le programme résout un p-médian non capacitaire.
Les variables $X_l\in\{0,\dots,m\}$ indiquent le nombre de médianes placées
au nœud $l$. La colocalisation est donc autorisée.

Le modèle impose :
\begin{align}
\sum_{l\in N}X_l &= m,\\
z^{P}_{lj} &\leq w_jX_l,
  &&l,j\in N,\\
\sum_{l\in N}z^{P}_{lj} &= w_j,
  &&j\in N.
\end{align}

Il maximise :
\[
\sum_{l,j}(D_{\max}-d_{lj})z^{P}_{lj},
\]
ce qui est équivalent à minimiser la distance totale puisque chaque demande
est ici imposée à l'égalité.

La solution $X$ est développée en $m$ établissements individuels. Pour
chaque établissement $i$, la capacité médiane $C^m_i$ est la somme des flux
affectés à sa médiane. Cette capacité sert ensuite de référence à tous les
scénarios stochastiques.

\section{Algorithmes de résolution}

\subsection{Évaluation de la solution p-médiane}

La localisation p-médiane est fixée dans le modèle stochastique et les flux
de tous les scénarios sont optimisés. La valeur obtenue est
\code{VMEDIAN1}.

Une recherche locale part ensuite de cette localisation ; sa valeur est
\code{VMEDIAN2}.

\subsection{Heuristique gloutonne}

La méthode active est \code{CMpuf::MyGreedyHeuristic}. Les établissements
sont placés séquentiellement :
\begin{enumerate}
  \item à l'étape $i$, les établissements $j>i$ sont temporairement forcés à
        zéro ;
  \item la contrainte \eqref{eq:localisation} de l'établissement $i$ est
        ajoutée ;
  \item CPLEX optimise le modèle partiel ;
  \item le meilleur emplacement de $i$ est mémorisé et fixé à un ;
  \item la contrainte temporaire est retirée et l'étape suivante commence.
\end{enumerate}

La valeur finale est \code{VGREEDY} et la localisation est notée $L^G$.

\subsection{Recherche locale}

La méthode \code{Localsearch} met en œuvre un meilleur déplacement à un
établissement :
\begin{enumerate}
  \item toutes les positions courantes sont fixées ;
  \item chaque établissement $i$ est libéré à tour de rôle ;
  \item le modèle est résolu pour trouver son meilleur nouvel emplacement ;
  \item le meilleur déplacement strictement améliorant est accepté ;
  \item le processus recommence jusqu'à ce qu'aucun déplacement ne soit
        améliorant.
\end{enumerate}

Si le rayon demandé est inférieur à $n$, l'établissement $i$ ne peut être
déplacé que vers les \code{rayon} nœuds les plus proches de sa position
initiale. Les nœuds sont triés par \code{ROPTri}.

La recherche est lancée deux fois :
\begin{itemize}
  \item depuis la solution gloutonne : \code{VGREEDYLS} ;
  \item depuis la solution p-médiane : \code{VMEDIAN2}.
\end{itemize}

\subsection{Résolution exacte}

Si le deuxième argument vaut \code{e}, le MIP complet est résolu par CPLEX.
Les paramètres actifs sont :
\begin{itemize}
  \item limite de temps : argument \code{LIMITE\_TEMPS} ;
  \item gap relatif : argument \code{GAP} ;
  \item fréquence des heuristiques internes : $-1$ ;
  \item stratégie Benders : automatique ($-1$).
\end{itemize}

Les statuts \emph{Optimal}, \emph{OptimalTol}, \emph{SolLim} et
\emph{AbortTimeLim} sont acceptés pour récupérer une solution. La valeur
\code{VOPT} peut donc être une meilleure solution connue à la limite de temps,
et pas nécessairement une valeur prouvée optimale.

Le meilleur résultat heuristique est calculé dans \code{bestL}, mais l'appel
\code{MIPStart} est actuellement commenté : ce point de départ n'est pas
transmis à CPLEX.

\subsection{Résolution indépendante des scénarios}

Après une résolution exacte, \code{ResolutionScenarios} résout un modèle de
transport/localisation séparé pour chaque scénario avec sa capacité connue.
Ces solutions servent à alimenter \code{SolScenarios} et \code{SolSupply}.
La fonction d'affichage détaillé associée existe, mais son appel est commenté
dans le programme principal.

\section{Indicateurs complémentaires}

Pour une localisation $L=(L_0,\dots,L_{m-1})$, la centralisation écrite par
le programme est la distance moyenne entre paires d'établissements :
\[
\operatorname{cent}(L)
=
\frac{\sum_{0\leq i<j<m}d_{L_iL_j}}
     {m(m-1)/2}.
\]

Si $q(L)$ est le nombre d'emplacements distincts occupés, la colocalisation
est :
\[
\operatorname{coloc}(L)
=
100\frac{m-q(L)}{m}.
\]
Elle vaut zéro si tous les établissements occupent des nœuds distincts.

\section{Formats des fichiers de test}

\subsection{Fichier liste d'instances}

Le premier argument de l'exécutable désigne un fichier texte contenant un
chemin d'instance par ligne :
\begin{lstlisting}
/chemin/absolu/vers/instance1
/chemin/absolu/vers/instance2
\end{lstlisting}

Par exemple, \code{../BENCHMARK/calgaryhospital15} contient le chemin de
l'instance réelle. Les chemins ne doivent pas contenir d'espace, car ils sont
lus avec \code{fscanf(file, "\%s", filename)}.

\subsection{Format standard ou hospital}

Lorsque \code{TYPE\_ENTREE} est différent de \code{geometric}, le fichier
d'instance est lu comme une suite de nombres :
\begin{lstlisting}
n
w[0] w[1] ... w[n-1]
d[0][0] ... d[0][n-1]
...
d[n-1][0] ... d[n-1][n-1]
\end{lstlisting}

Les espaces, tabulations et retours à la ligne sont équivalents pour
l'opérateur de lecture C++.

\subsection{Format geometric}

Lorsque le dernier argument vaut exactement \code{geometric}, le lecteur
attend :
\begin{lstlisting}
n
mtmp
ptmp[0] ... ptmp[mtmp-1]
w[0] ... w[n-1]
d[0][0] ... d[n-1][n-1]
\end{lstlisting}

Le tableau temporaire \code{ptmp} est lu mais n'est pas utilisé dans les
calculs ultérieurs. Le vecteur des demandes et la matrice sont lus au moyen
des opérateurs Concert sur \code{IloNumArray} et \code{IloArray}.

\subsection{Fichier de paramètres stochastiques}

Le fichier \code{parameters/param} respecte :
\begin{lstlisting}
ncap
a[0] ... a[ncap-2]
b[0] ... b[ncap-2]
\end{lstlisting}

Le fichier fourni est :
\begin{lstlisting}
4
0.10 0.40 0.50
0.00 0.33 0.66
\end{lstlisting}

\section{Procédure de compilation et de test}

\subsection{Compilation}

Depuis la racine :
\begin{lstlisting}[language=bash]
make clean
make -j2 cmpuf
\end{lstlisting}

ou avec CMake :
\begin{lstlisting}[language=bash]
cmake --preset release
cmake --build --preset cmpuf --parallel
\end{lstlisting}

L'exécutable produit est \code{bin/cmpuf}.

\subsection{Interface en ligne de commande}

Le programme attend neuf arguments utilisateur :
\begin{lstlisting}[language=bash]
./bin/cmpuf LISTE_INSTANCES MODE M GAP RAYON_LOCAL \
  RAYON_BB LIMITE_TEMPS FICHIER_PARAMETRES TYPE_ENTREE
\end{lstlisting}

\begin{longtable}{clp{8.8cm}}
\toprule
No & Argument & Rôle\\
\midrule
\endhead
1 & \code{LISTE\_INSTANCES} & Fichier contenant un chemin d'instance par
ligne.\\
2 & \code{MODE} & \code{e} active la résolution exacte ; toute autre valeur
conserve les évaluations heuristiques.\\
3 & \code{M} & Nombre d'établissements.\\
4 & \code{GAP} & Gap relatif CPLEX.\\
5 & \code{RAYON\_LOCAL} & Taille du voisinage spatial de recherche locale.\\
6 & \code{RAYON\_BB} & Rayon prévu pour le branch-and-bound local ; les appels
correspondants sont actuellement commentés.\\
7 & \code{LIMITE\_TEMPS} & Limite de temps CPLEX en secondes.\\
8 & \code{FICHIER\_PARAMETRES} & Paramètres $ncap$, $a_r$ et $b_r$.\\
9 & \code{TYPE\_ENTREE} & \code{geometric} ou format standard, par convention
\code{hospital}.\\
\bottomrule
\end{longtable}

Exemple :
\begin{lstlisting}[language=bash]
./bin/cmpuf ../BENCHMARK/calgaryhospital15 e 4 0.00 \
  5 40 30000 parameters/param hospital
\end{lstlisting}

\subsection{Plan expérimental codé}

Pour chaque instance de la liste :
\begin{enumerate}
  \item le p-médian est résolu une fois ;
  \item $\gamma$ parcourt $1$, $1{,}25$, $1{,}50$, $1{,}75$ ;
  \item pour chaque $\gamma$, $\alpha$ parcourt $0$, $0{,}1$, \ldots, $1$ ;
  \item $\beta$ reste fixé à zéro ;
  \item pour chaque couple $(\alpha,\gamma)$, le programme évalue le
        p-médian, le glouton et les deux recherches locales ;
  \item en mode \code{e}, il lance aussi le MIP exact.
\end{enumerate}

Cela représente théoriquement $4\times 11=44$ configurations par instance,
sous réserve des effets d'arrondi de la boucle en nombres flottants.

\subsection{Contrôles recommandés}

\begin{enumerate}
  \item lancer \code{./bin/cmpuf} sans argument et vérifier le message
        \og Erreur sur le nombre d'arguments \fg ;
  \item commencer avec une petite instance et une limite de temps courte ;
  \item vérifier que $\sum_{r<ncap-1}a_r=1$ ;
  \item vérifier que $m$, $ncap^m$ et $n^2ncap^m$ restent compatibles avec la
        mémoire disponible ;
  \item contrôler que le répertoire \code{results/hospital} existe et est
        accessible en écriture ;
  \item comparer \code{VGREEDY}, \code{VGREEDYLS}, \code{VMEDIAN1},
        \code{VMEDIAN2} et \code{VOPT}.
\end{enumerate}

\section{Formats des fichiers de résultats}

\subsection{Construction des noms}

Quatre préfixes sont codés en dur dans \code{cpp/cmpuf.cpp} :
\begin{lstlisting}
results2026_v1_
solgreedy_v1_
solopt_v1_
solcent_v1_
\end{lstlisting}

Le suffixe concatène une partie de \code{argv[1]}, puis :
\[
\code{\_M\_GAP\_RAYON\_LOCAL\_RAYON\_BB\_LIMITE\_TEMPS}.
\]

\paragraph{Limite actuelle.}
La fonction \code{concat} utilise \code{argv[1]+16} au lieu d'extraire
proprement le dernier composant du chemin. Le nom produit dépend donc de la
longueur du chemin de la liste d'instances et peut être tronqué, comme
\code{garyhospital15} au lieu de \code{calgaryhospital15}. Les répertoires de
sortie sont également des chemins absolus codés en dur.

Tous les fichiers sont ouverts en mode ajout (\code{ios::app} ou
\code{fopen(...,"a")}).

\subsection{Fichier principal results2026}

Chaque ligne de données contient 25 champs séparés par \code{\&} :
\begin{longtable}{clp{8.5cm}}
\toprule
No & Champ & Signification\\
\midrule
\endhead
1 & \code{n} & Nombre de clients/nœuds.\\
2 & \code{m} & Nombre d'établissements.\\
3--5 & \code{Alpha}, \code{Beta}, \code{Gamma} & Paramètres expérimentaux.\\
6 & \code{VOPT} & Valeur du MIP exact, ou $-1$ hors mode exact.\\
7 & \code{VGREEDY} & Valeur gloutonne.\\
8 & \code{VGREEDYLS} & Valeur après recherche locale depuis le glouton.\\
9 & \code{VMEDIAN1} & Valeur stochastique de la localisation p-médiane.\\
10 & \code{VMEDIAN2} & Valeur après recherche locale depuis le p-médian.\\
11--15 & \code{cent\_*} & Centralisations des cinq solutions.\\
16--20 & \code{coloc\_*} & Pourcentages de colocalisation.\\
21 & \code{topt} & Temps CPLEX de la résolution exacte.\\
22 & \code{tgreedy} & Temps de l'heuristique gloutonne.\\
23 & \code{tgreedyls} & Temps de recherche locale depuis le glouton.\\
24 & \code{tpmedian1} & Temps du prétraitement p-médiane.\\
25 & \code{tpmedian2} & Temps de recherche locale depuis le p-médian.\\
\bottomrule
\end{longtable}

\paragraph{Anomalie actuelle de l'en-tête.}
L'appel qui écrit les libellés contient 24 marqueurs \code{\%s} pour 25
libellés, et n'ajoute pas de retour à la ligne. Le libellé
\code{tpmedian2} n'est donc pas écrit et la première ligne de données est
concaténée à l'en-tête. Les lignes numériques suivantes respectent néanmoins
les 25 champs du format \code{format1}.

\subsection{Fichier solgreedy}

Une ligne est ajoutée pour chaque couple $(\alpha,\gamma)$. Elle contient les
$m$ indices de nœuds de la solution obtenue après recherche locale depuis le
glouton :
\begin{lstlisting}
Lls1[0] Lls1[1] ... Lls1[m-1]
\end{lstlisting}
Malgré le préfixe \code{solgreedy}, il s'agit donc de $L^{G,LS}$, pas de la
localisation gloutonne brute.

\subsection{Fichier solopt}

En mode exact, une ligne contient les $m$ indices de la solution récupérée
après le MIP :
\begin{lstlisting}
Lopt[0] Lopt[1] ... Lopt[m-1]
\end{lstlisting}
Aucune ligne n'est ajoutée hors mode \code{e}.

\subsection{Fichier solcent}

En mode exact, chaque ligne comporte cinq colonnes séparées par des
tabulations :
\begin{lstlisting}
alpha    beta    gamma    VOPT    centralizationopt
\end{lstlisting}
Les nombres sont écrits en notation décimale fixe.

\subsection{Sortie standard}

Pour chaque configuration, le programme affiche une ligne de style tableau
LaTeX contenant d'abord le dernier composant de \code{filename}, puis les 25
valeurs décrites ci-dessus. Cette sortie peut être redirigée :
\begin{lstlisting}[language=bash]
./bin/cmpuf ... > compte_rendu.txt
\end{lstlisting}

\section{Points de vigilance pour l'interprétation}

\begin{itemize}
  \item L'objectif est une utilité attendue, pas directement le coût de
        transport ; la conversion est exacte seulement si les quantités
        totales livrées sont fixées.
  \item Les demandes stochastiques sont imposées par $\geq$ ; un sur-service
        est mathématiquement possible.
  \item La colocalisation est autorisée dans le p-médian et dans le modèle
        stochastique actif.
  \item \code{beta} est actuellement sans effet sur les capacités.
  \item Le nombre de scénarios et de variables croît exponentiellement avec
        $m$.
  \item Une valeur \code{VOPT} obtenue à la limite de temps n'est pas
        nécessairement prouvée optimale.
  \item Les fichiers sont complétés en mode ajout ; deux campagnes de même nom
        sont concaténées.
  \item La boucle \code{while(!feof(file1))} peut tenter une itération
        supplémentaire ; les listes d'instances doivent être soigneusement
        terminées et les sorties contrôlées.
\end{itemize}

\section{Traçabilité entre formulation et code}

\begin{longtable}{p{5.2cm}p{9cm}}
\toprule
Élément & Fonction ou fichier principal\\
\midrule
\endhead
Lecture des instances & Constructeur \code{CMpuf}, \code{cpp/cmpuf.cpp}\\
Lecture des paramètres & Constructeur \code{CMpuf}, \code{cpp/cmpuf.cpp}\\
Capacités & \code{CMpuf::Init\_cap}\\
Scénarios et probabilités & \code{CMpuf::Init\_T}, \code{CMpuf::Init\_C}\\
Variables $x,z$ & \code{CMpuf::Init\_x}, \code{CMpuf::Init\_z}\\
Contraintes & \code{CMpuf::Init\_contraintes},
\code{CMpuf::Init\_contraintes\_x}\\
Objectif & \code{CMpuf::Init\_obj}\\
P-médian & \code{cpp/pmedian.cpp}\\
Glouton & \code{CMpuf::MyGreedyHeuristic}\\
Recherche locale & \code{CMpuf::Localsearch}\\
Résolution exacte & \code{CMpuf::Resolution} et programme principal\\
Scénarios isolés & \code{CMpuf::ResolutionScenarios},
\code{cpp/transport.cpp}\\
Résultats & \code{AffichageSolutions},
\code{AffichageSolutionsScenarios}, \code{AffichageCoutTransport}\\
\bottomrule
\end{longtable}

\end{document}
